% GENERATED FILE. Edit canonical lessons, then run npm run generate:books.
% canonical-source-hash: fnv1a64:276e29b2a650625d
% canonical-lessons: FR-C370-imprevu, FR-C370-previsible, FR-C370-imprevisible, FR-C370-adapter, FR-C370-souple

\chapter{Les imprévus}
\label{ch:fr-les-imprevus}
\hlchaptermodality{370}

\begin{chapteropening}
\textbf{By the end of this chapter:} \emph{I can talk about the unexpected: a hitch, what is foreseeable or not, adapting to it and being easygoing.}

Complete the last lesson of chapter 370: I can talk about the unexpected: a hitch, what is foreseeable or not, adapting to it and being easygoing.
\end{chapteropening}

\section[l'imprévu]{l'imprévu — the hitch, the unexpected event}

\label{lesson:FR-C370-imprevu}



\begin{quote}
Before the new one: say the French for the possibility, then the French for for lack of.
\end{quote}

\subsection*{You'll want to know: l'imprévu}
\textbf{l'imprévu} — \textquotedblleft{}the hitch, the unexpected event\textquotedblright{}.

\begin{quote}
J'ai un imprévu, je ne peux pas venir ce soir.
\end{quote}

\begin{quote}
\emph{Something has come up, I cannot come tonight.}
\end{quote}

Also an adjective: \emph{une visite imprévue}. Masculine.

\begin{cousinweb}
\emph{imprévu} is built from \emph{in-} \textquotedblleft{}not\textquotedblright{} and \emph{prévu}, the participle of \emph{prévoir} \textquotedblleft{}to foresee\textquotedblright{} (EtymDB), from Latin \emph{praevidēre} \textquotedblleft{}to foresee\textquotedblright{}, built on \emph{vidēre} \textquotedblleft{}to see\textquotedblright{} (Lewis \& Short).
\end{cousinweb}

\begin{grammarlens}[title={What you've built}]
One more word for telling, explaining and describing at B1.
\end{grammarlens}

\subsection*{Your turn}
\begin{itemize}
\raggedright
  \item \emph{Say it:} \emph{l'imprévu}
  \item \emph{Say it:} \emph{l'imprévu}, once more
  \item \emph{Say it:} a sentence of your own with \emph{l'imprévu}
  \item \emph{Recall:} say the French for the possibility, then the French for for lack of, then say \emph{l'imprévu} again
\end{itemize}

\begin{tcolorbox}[breakable,colback=teal!4,colframe=teal!35!black,title={Before you move on}]
What does l'imprévu mean? (\textquotedblleft{}The hitch, the unexpected event\textquotedblright{}.) Say it once more.
\end{tcolorbox}
\section[prévisible]{prévisible — foreseeable, predictable}

\label{lesson:FR-C370-previsible}



\begin{quote}
Before the new one: say the French for for lack of, then the French for the hitch.
\end{quote}

\subsection*{You'll want to know: prévisible}
\textbf{prévisible} — \textquotedblleft{}foreseeable, predictable\textquotedblright{}.

\begin{quote}
Ce retard était prévisible avec cette neige.
\end{quote}

\begin{quote}
\emph{This delay was predictable with this snow.}
\end{quote}

One form for both genders.

\begin{cousinweb}
\emph{prévisible} is built on \emph{prévoir} \textquotedblleft{}to foresee\textquotedblright{} and the suffix \emph{-ible}. \emph{prévoir} comes from Latin \emph{praevidēre} \textquotedblleft{}to foresee\textquotedblright{}, built on \emph{vidēre} \textquotedblleft{}to see\textquotedblright{} (Lewis \& Short).
\end{cousinweb}

\begin{grammarlens}[title={What you've built}]
One more word for telling, explaining and describing at B1.
\end{grammarlens}

\subsection*{Your turn}
\begin{itemize}
\raggedright
  \item \emph{Say it:} \emph{prévisible}
  \item \emph{Say it:} \emph{prévisible}, once more
  \item \emph{Say it:} a sentence of your own with \emph{prévisible}
  \item \emph{Recall:} say the French for for lack of, then the French for the hitch, then say \emph{prévisible} again
\end{itemize}

\begin{tcolorbox}[breakable,colback=teal!4,colframe=teal!35!black,title={Before you move on}]
What does prévisible mean? (\textquotedblleft{}Foreseeable, predictable\textquotedblright{}.) Say it once more.
\end{tcolorbox}
\section[imprévisible]{imprévisible — unpredictable}

\label{lesson:FR-C370-imprevisible}



\begin{quote}
Before the new one: say the French for the hitch, then the French for foreseeable.
\end{quote}

\subsection*{You'll want to know: imprévisible}
\textbf{imprévisible} — \textquotedblleft{}unpredictable\textquotedblright{}.

\begin{quote}
Le temps à la montagne est imprévisible.
\end{quote}

\begin{quote}
\emph{The weather in the mountains is unpredictable.}
\end{quote}

With \emph{in-}, the opposite of \emph{prévisible}.

\begin{cousinweb}
\emph{imprévisible} is built from \emph{in-} \textquotedblleft{}not\textquotedblright{} and \emph{prévisible}, which rests on \emph{prévoir}, from Latin \emph{praevidēre} \textquotedblleft{}to foresee\textquotedblright{} (Lewis \& Short).
\end{cousinweb}

\begin{grammarlens}[title={What you've built}]
One more word for telling, explaining and describing at B1.
\end{grammarlens}

\subsection*{Your turn}
\begin{itemize}
\raggedright
  \item \emph{Say it:} \emph{imprévisible}
  \item \emph{Say it:} \emph{imprévisible}, once more
  \item \emph{Say it:} a sentence of your own with \emph{imprévisible}
  \item \emph{Recall:} say the French for the hitch, then the French for foreseeable, then say \emph{imprévisible} again
\end{itemize}

\begin{tcolorbox}[breakable,colback=teal!4,colframe=teal!35!black,title={Before you move on}]
What does imprévisible mean? (\textquotedblleft{}Unpredictable\textquotedblright{}.) Say it once more.
\end{tcolorbox}
\section[adapter]{adapter — to adapt (mostly reflexive, s'adapter)}

\label{lesson:FR-C370-adapter}



\begin{quote}
Before the new one: say the French for foreseeable, then the French for unpredictable.
\end{quote}

\subsection*{You'll want to know: adapter}
\textbf{adapter} — \textquotedblleft{}to adapt (mostly reflexive, s'adapter)\textquotedblright{}.

\begin{quote}
Les enfants s'adaptent vite à une nouvelle école.
\end{quote}

\begin{quote}
\emph{Children adapt quickly to a new school.}
\end{quote}

Mostly reflexive: \emph{s'adapter à} — to adapt to. \emph{Adapter un livre} — to adapt a book for the screen.

\begin{cousinweb}
\emph{adapter} is a learned borrowing from Latin \emph{adaptāre} \textquotedblleft{}to fit to\textquotedblright{} (EtymDB), built from \emph{ad} \textquotedblleft{}to\textquotedblright{} and \emph{aptāre} \textquotedblleft{}to fit\textquotedblright{}, from \emph{aptus} \textquotedblleft{}fitted\textquotedblright{} (Lewis \& Short).
\end{cousinweb}

\begin{grammarlens}[title={What you've built}]
One more word for telling, explaining and describing at B1.
\end{grammarlens}

\subsection*{Your turn}
\begin{itemize}
\raggedright
  \item \emph{Say it:} \emph{adapter}
  \item \emph{Say it:} \emph{adapter}, once more
  \item \emph{Say it:} a sentence of your own with \emph{adapter}
  \item \emph{Recall:} say the French for foreseeable, then the French for unpredictable, then say \emph{adapter} again
\end{itemize}

\begin{tcolorbox}[breakable,colback=teal!4,colframe=teal!35!black,title={Before you move on}]
What does adapter mean? (\textquotedblleft{}To adapt (mostly reflexive, s'adapter)\textquotedblright{}.) Say it once more.
\end{tcolorbox}
\section[souple]{souple — supple; easygoing, adaptable}

\label{lesson:FR-C370-souple}



\begin{quote}
Before the new one: say the French for unpredictable, then the French for to adapt.
\end{quote}

\subsection*{You'll want to know: souple}
\textbf{souple} — \textquotedblleft{}supple; easygoing, adaptable\textquotedblright{}.

\begin{quote}
Mes horaires sont souples : je peux venir plus tard.
\end{quote}

\begin{quote}
\emph{My hours are adaptable: I can come later.}
\end{quote}

Of the body: supple. Of a person or a rule: easygoing.

\begin{cousinweb}
\emph{souple} is inherited from Latin \emph{supplex} \textquotedblleft{}bending the knee, humble\textquotedblright{} (EtymDB; Lewis \& Short), built on \emph{plicāre} \textquotedblleft{}to fold\textquotedblright{}.
\end{cousinweb}

\begin{grammarlens}[title={What you've built}]
One more word for telling, explaining and describing at B1.
\end{grammarlens}

\subsection*{Your turn}
\begin{itemize}
\raggedright
  \item \emph{Say it:} \emph{souple}
  \item \emph{Say it:} \emph{souple}, once more
  \item \emph{Say it:} a sentence of your own with \emph{souple}
  \item \emph{Recall:} say the French for unpredictable, then the French for to adapt, then say \emph{souple} again
\end{itemize}

\begin{tcolorbox}[breakable,colback=teal!4,colframe=teal!35!black,title={Before you move on}]
What does souple mean? (\textquotedblleft{}Supple; easygoing, adaptable\textquotedblright{}.) Say it once more.
\end{tcolorbox}
